\subsection{The Binomial Probability Distribution}
\hrulefill

\paragraph*{Definition:}
There are many experiments that conform either exactly or approximately to the following list of requirements:
\begin{itemize}
    \item The experiment consists of a sequence of $n$ identical trials.
    \item Each trial results in one of two outcomes, which we will call a success and a failure.
    \item The trials are independent, so that the outcome on any particular trial does not influence the outcome on any other trial.
    \item The probability of a success, denoted by $p$, remains the same from trial to trial. The probability of a failure is $q = 1 - p$.
\end{itemize}
An experiment that satisfies these requirements is called a \textbf{binomial experiment}.

\paragraph*{Example:}
Consider each of the next $n$ vehicles undergoing an emissions test, and let S denote a cehicle thta passes the test and $F$ denote one that fails.
\begin{itemize}
    \item The experiment consists of $n$ identical trials, where each trial is the emissions test of a vehicle.
    \item Each trial results in one of two outcomes: S or F.
    \item The trials are independent, assuming that the vehicles are randomly selected.
    \item The probability of a success (a vehicle passing the test) remains the same from trial to trial, assuming that the vehicles are randomly selected.
\end{itemize}
Thus, this is a binomial experiment.

\subsubsection*{The Binomail Random Variable and Distribution}

In most binomial experiments, we are interested in the total number of successes that occur in the $n$ trials rather than the specific sequence of
 successes and failures.
\paragraph*{Definition:}
The \textbf{binomial random variable} $X$ associated with a binomial experiment consisting of $n$ trials is defined as:
\[
    X = \text{the number of successes in } n \text{ trials.}
\]

\paragraph*{Example:}
Suppose that $n=3$. There are 8 possible outcomes of the experiment, as shown in the following table:
\begin{center}
    \begin{tabular}{c|c}
        Outcome & $X$ \\
        \hline
        SSS & 3 \\
        SSF & 2 \\
        SFS & 2 \\
        FSS & 2 \\
        SFF & 1 \\
        FSF & 1 \\
        FFS & 1 \\
        FFF & 0 \\
    \end{tabular}
\end{center}
The possible values of $X$ are 0, 1, 2, and 3.

\paragraph*{Example:}
Consider the case that $n=4$. Find P(SSFS) and P(SSSS):
\begin{align*}
    P(SSFS) &= P(S) \cdot P(S) \cdot P(F) \cdot P(S) \\
    &= p \cdot p \cdot (1-p) \cdot p \\
    &= p^3(1-p)
    \\
    P(SSSS) &= P(S) \cdot P(S) \cdot P(S) \cdot P(S) \\
    &= p \cdot p \cdot p \cdot p \\
    &= p^4
\end{align*}

\pagebreak

\paragraph*{Notation:}
Because the pmf of a binomial random variable $X$ depends on the two parameters $n$ and $p$, we denote the pmf by $b(x: n, p)$.

\paragraph*{Example:}
Find $b(3; 4, p)$:
\begin{align*}
    b(3; 4, p) &= P(X=3) \\
    &= P(SSSF) + P(SSFS) + P(SFSS) + P(FSSS) \\
    &= p^3(1-p) + p^3(1-p) + p^3(1-p) + p^3(1-p) \\
    &= 4p^3(1-p)
\end{align*}

\paragraph{Theorem}:
\[
    b(x; n, p) = \begin{cases}
        \binom{n}{x} p^x (1-p)^{n-x}, & x = 0, 1, 2, \ldots, n \\
        0, & \text{otherwise}
    \end{cases}
\]

\paragraph*{Example:}
Each of six randomly selected cola drinkers is given a glass containing cola $S$ and one containing cola $F$. 
The glasses are identical in appearance except for a code on the bottom to identify the cola. 
Suppose there is actually no tendency among cola drinkers to prefer one cola over the other. 
Then $p = P(S) = 0.5$. Find (1) the probability that exactly four of the six cola drinkers prefer cola $S$,
(2) the probability that at least four of the six cola drinkers prefer cola $S$, and (3) the probability that at most
1 of the six cola drinkers prefer cola $S$.

\paragraph*{1.}
\begin{align*}
    P(X=4) = b(4; 6, 0.5) &= \binom{6}{4} (0.5)^4 (1-0.5)^{6-4} \\
    &= \binom{6}{4} (0.5)^4 (0.5)^2 \\
    &= 15 \cdot (0.5)^6 \\
    &= \frac{15}{64} \approx 0.2344
\end{align*}

\paragraph*{2.}
\begin{align*}
    P(X \geq 4) &= \sum_{x=4}^{6} b(x; 6, 0.5) \\
    &= \sum_{x=4}^{6} \binom{6}{x} (0.5)^x (1-0.5)^{6-x} \\
    &= \frac{22}{64} \approx 0.3438
\end{align*}

\paragraph*{3.}
\begin{align*}
    P(X \leq 1) &= \sum_{x=0}^{1} b(x; 6, 0.5) \\
    &= \sum_{x=0}^{1} \binom{6}{x} (0.5)^x (1-0.5)^{6-x} \\
    &= \frac{7}{64} \approx 0.1094
\end{align*}

\subsubsection*{Binomial Tables}
For $X \sim \text{Binomial}(n, p)$, the cdf will be denoted by:
\[
    B(x; n, p) = P(X \leq x) = \sum_{y=0}^{\lfloor x \rfloor} b(y; n, p)
\]

\paragraph*{Theorem:}
If $X \sim B(n,p)$, then:
\begin{itemize}
    \item $E(X) = np$
    \item $V(X) = np(1-p)$
    \item $\sigma_X = \sqrt{np(1-p)}$
\end{itemize}
